\documentclass[11pt]{article}
\usepackage[margin=1.1in]{geometry}
\usepackage{amsmath, amssymb, amsthm, mathtools}
\usepackage{booktabs}
\usepackage{enumitem}
\usepackage{hyperref}
\usepackage{parskip}
\usepackage{xcolor}
\hypersetup{colorlinks=true, urlcolor=blue, linkcolor=black}

\title{Comparing children to LMs on a single sigmoid-slope axis\\
\large{A derivation of the Chang \& Bergen bridge}}
\author{Standard Model 2 working note}
\date{\today}

\newcommand{\E}{\mathbb{E}}
\newcommand{\Var}{\mathrm{Var}}
\newcommand{\logit}{\mathrm{logit}}

\begin{document}
\maketitle

\section{The claim}

Children's per-word acquisition curves and large language models'
per-word acquisition curves both have a sigmoidal form in
$\log(\text{experience})$. They live on the same axis. The
\emph{slope} of that sigmoid is the natural object to compare
across the two populations.

This note derives the slope from each side, lines them up in matched
units, and reports the headline numbers we use in the deck.

\section{Children's side: per-child sigmoid slope is $\kappa_i$}

\subsection{Linear predictor}

In our model (slides 9--12), for child $i$, word $j$, age $t$:
\begin{equation}
\eta_{ij}(t) = \log r_i + \log\alpha_i + \kappa_i\,\log t - \delta_j,
\qquad
P(\text{produces}_{ij} \mid t) = \sigma(\eta_{ij}(t)),
\label{eq:our-model}
\end{equation}
where $\sigma(x) = 1/(1+e^{-x})$ is the standard logistic.
Per-child parameters: input rate $r_i$, efficiency $\alpha_i$,
scaling exponent $\kappa_i$. Per-word: difficulty $\delta_j$.

\subsection{Per-child slope}

Differentiate $\eta_{ij}$ with respect to $\log t$, holding
$(i, j)$ fixed:
\begin{equation}
\frac{\partial \eta_{ij}}{\partial \log t} = \kappa_i.
\label{eq:per-child-slope}
\end{equation}
Therefore $\logit P(\text{produces}_{ij}\mid t)$ is a
\emph{linear} function of $\log t$ with slope exactly $\kappa_i$.
The per-child sigmoid slope on $\log t$ is $\kappa_i$ --- by
construction.

\subsection{Population-marginalized slope}

If instead we ask, ``what slope does a researcher fit to
\emph{cross-sectional} Wordbank data for word $w$, marginalizing
over the population of kids?'', the answer attenuates $\kappa_{\text{pop}}$.

For fixed word $w$ and time $t$, define
\begin{equation*}
X_i = \log r_i + \log\alpha_i + \kappa_i \log t,
\end{equation*}
so $P(\text{produces}_{iw}\mid t) = \sigma(X_i - \delta_w)$.
With $(\log r_i, \log\alpha_i, \kappa_i)$ joint normal and (for
simplicity) treating cross-correlations as zero,
\begin{align}
\E[X_i] &= \mu_r + \kappa_{\text{pop}}\,\log t, \\
\Var[X_i] &= \sigma_r^2 + \sigma_\alpha^2 + \sigma_\kappa^2\,\log^2 t
        \;=\; \sigma_\xi^2 + \sigma_\kappa^2\,\log^2 t \;=:\; \sigma_X^2(t).
\end{align}

The cross-sectional acquisition probability is
\begin{equation*}
P_{\text{pop}}(t, w) = \E_i\!\left[\sigma(X_i - \delta_w)\right].
\end{equation*}
Using the standard logistic--Gaussian convolution approximation with
$a^2 = \pi^2/8 \approx 1.234$ (Daunizeau, 2017),
\begin{equation}
P_{\text{pop}}(t, w) \;\approx\;
\sigma\!\left(
  \frac{\mu_r + \kappa_{\text{pop}} \log t - \delta_w}
       {\sqrt{1 + a^2\,\sigma_X^2(t)}}
\right).
\label{eq:pop-marginal}
\end{equation}

Differentiating with respect to $\log t$ at the midpoint
$\log t_{50\%}$ where the numerator vanishes:
\begin{equation}
\boxed{
\left.\frac{\partial \logit P_{\text{pop}}}{\partial \log t}\right|_{t_{50\%}}
\;=\;
\frac{\kappa_{\text{pop}}}{\sqrt{1 + a^2\,\sigma_X^2(t_{50\%})}}.
}
\label{eq:pop-slope}
\end{equation}

For our M\_best on English: $\kappa_{\text{pop}} = 10.4$, $\sigma_\xi^2 = 1.80^2 \approx 3.24$, $\sigma_\kappa^2 = 3.47^2 \approx 12.0$. Plugging in:
\begin{itemize}\setlength\itemsep{0.2em}
\item At $t_{50\%}=a_0$ ($\log t = 0$): $\sigma_X^2 = 3.24$, slope $\approx 10.4/\sqrt{1+1.234\cdot 3.24} = 10.4/2.24 \approx 4.6$.
\item At $\log t = 0.4$ (about 28~mo if $a_0=19$): $\sigma_X^2 \approx 3.24 + 0.16\cdot 12 = 5.16$, slope $\approx 10.4/\sqrt{7.37} \approx 3.83$.
\end{itemize}

So a researcher fitting a single sigmoid to Wordbank cross-sectional
production data per word would observe \emph{population-marginalized}
slopes of about $4$~logits per natural-log-time unit. The per-child
slope is much higher ($\kappa_i \sim 10.4 \pm 3.5$); the cross-sectional
fit averages over kids and the slope shrinks.

\subsection{Which slope is the right LLM comparand?}

Both kid-side slopes have meaningful comparands on the LM side. Chang
\& Bergen fit \emph{per-LM, per-word} sigmoids: there's no
``population of kids'' equivalent because each LM is a fixed
deterministic learner. The LM analog of $\kappa_i$ is therefore the
fitted slope $1/\text{scale}_w$ for one (LM, word) pair --- same
``instance'' (a single deterministic learner on a single word).
That's the comparison we make in the deck. The
population-marginalized slope is what one would compare to a
cross-LM-architecture pooled fit.

\section{LM side: Chang \& Bergen sigmoid fit}

\subsection{Their parameterization}

For each (LM, word) pair, Chang \& Bergen (2022) fit per-word mean
surprisal $S_w$ during training as a function of $x = \log_{10}(\text{steps})$,
using the four-parameter logistic
\begin{equation}
S_w(x) = \mathrm{ParamLower}_w + \frac{\mathrm{ParamUpper}_w - \mathrm{ParamLower}_w}
                                       {1 + \exp((x - \mathrm{ParamXmid}_w)/\mathrm{ParamScale}_w)}.
\label{eq:cb-fit}
\end{equation}
$\mathrm{ParamUpper}_w$ is the asymptote at $x \to -\infty$ (high
surprisal at training start), $\mathrm{ParamLower}_w$ is the asymptote
at $x \to +\infty$ (low surprisal after learning).
$\mathrm{ParamXmid}_w$ is the $x$ at the curve's midpoint, and
$\mathrm{ParamScale}_w$ is the scale parameter.

\subsection{Translating to ``proportion learned''}

Define proportion-learned as the fractional surprisal reduction:
\begin{equation}
P_{w}^{\mathrm{LM}}(x)
\;:=\;
\frac{\mathrm{ParamUpper}_w - S_w(x)}
     {\mathrm{ParamUpper}_w - \mathrm{ParamLower}_w}
\;=\;
1 \;-\; \frac{1}{1+\exp((x - \mathrm{ParamXmid}_w)/\mathrm{ParamScale}_w)}
\;=\;
\sigma\!\!\left(\frac{x - \mathrm{ParamXmid}_w}{\mathrm{ParamScale}_w}\right).
\end{equation}
So $P_w^{\mathrm{LM}}$ is a sigmoid in $x = \log_{10}(\text{steps})$
with slope on $\logit$:
\begin{equation}
\frac{\partial \logit P_w^{\mathrm{LM}}}{\partial x} \;=\; \frac{1}{\mathrm{ParamScale}_w}.
\end{equation}

\subsection{Unit conversion from $\log_{10}$ to natural log}

Let $y = \ln(\text{steps})$. Since $x = y/\ln 10$,
\begin{equation}
\boxed{
\frac{\partial \logit P_w^{\mathrm{LM}}}{\partial y}
\;=\;
\frac{1}{\mathrm{ParamScale}_w}\cdot\frac{1}{\ln 10}
\;\approx\;
\frac{0.434}{\mathrm{ParamScale}_w}.
}
\label{eq:lm-slope}
\end{equation}
Per-word, per-LM sigmoid slope on \emph{natural-log} training experience.

\section{Direct comparison}

Equations~\eqref{eq:per-child-slope} and \eqref{eq:lm-slope} are
on identical units: ``logit per natural-log unit of experience.''
Reading from our M\_best fits and Chang \& Bergen's processed sigmoid
fits:

\begin{center}
\renewcommand{\arraystretch}{1.20}
\begin{tabular}{@{}lcc@{}}
\toprule
& \textbf{median slope} & \textbf{IQR} \\
\midrule
Children, per-child $\kappa_i$ (English M\_best) & $10.3$ & $[8.0,\,12.6]$ \\
Children, per-child $\kappa_i$ (Norwegian M\_best) & $12.5$ & $[10.0,\,14.9]$ \\
LMs, BERT per-word                                & $0.76$ & $[0.42,\,1.32]$ \\
LMs, BiLSTM per-word                              & $0.87$ & $[0.53,\,1.52]$ \\
LMs, GPT-2 per-word                               & $0.81$ & $[0.45,\,1.54]$ \\
LMs, LSTM per-word                                & $0.96$ & $[0.44,\,1.90]$ \\
\bottomrule
\end{tabular}
\end{center}

\textbf{Headline.} Children are 10--15$\times$ steeper on the same
sigmoid-slope axis. LMs cluster at $\kappa \approx 1$, the unit
accumulator boundary; children sit in a wide distribution around
$\kappa \approx 10$.

\section{Caveats}

\begin{enumerate}[leftmargin=1.4em]
\item \textbf{Per-instance vs.\ per-instance, but different ``instance.''}
Kids' slopes vary across \emph{children}; LM slopes vary across
\emph{words} (within an LM, each word gets its own $1/\text{scale}_w$).
A single LM is a deterministic learner, so there's no
population-of-LM-instances analog of $\sigma_\kappa$ within architecture
--- variation is across the word set, not across replicates.

\item \textbf{Different experience units.} The $\log t$ axis on the kid
side is log-months-of-life. The $x$ axis on the LM side is
log-training-steps; for batch size 128 and average sentence length
$\sim$25, that's roughly log-(training-tokens/3200). On the log axis,
absolute alignment doesn't matter for slope; relative changes in
experience do.

\item \textbf{Cross-sectional Wordbank fit would give a smaller slope.}
A researcher who fit a sigmoid to Wordbank's per-word
production-vs-age data (one curve per word, marginalizing over kids)
would observe $\sim$4 on the same axis, not 10. That's the
population-marginalized slope from
equation~\eqref{eq:pop-slope}. The 10$\times$ headline is the
\emph{per-child} slope.

\item \textbf{LM training distribution is not CDS.} Chang \& Bergen
trained on BookCorpus + WikiText, which is adult written text, not
child-directed speech. Per-word acquisition slopes there reflect
distributional statistics of that corpus, not what kids hear. See
the next section.
\end{enumerate}

\section{Forward direction: CHILDES-trained GPT (Feng et al.\ 2026)}

Chang \& Bergen's comparison is a strong start but conflates two
factors: (1) the structural difference in scaling between humans and
LMs, and (2) the input distribution mismatch (adult-written-text vs.\
CDS). Feng et al.~(2026) train GPT-style models on CHILDES, which
matches the input distribution kids actually receive.

If their models retain training-step-by-step surprisal (or can be
rerun with that logging), we can:
\begin{enumerate}[leftmargin=1.4em]
\item Fit Chang \& Bergen-style sigmoids to CHILDES-trained GPT
surprisal trajectories on \emph{exactly the CDI vocabulary}.
\item Re-extract per-word slopes $1/\mathrm{scale}_w$ on natural-log
experience.
\item Replot Figure~18 of the deck (children-vs-LMs sigmoid-slope
distributions) with these matched-input slopes.
\end{enumerate}

The structural-difference hypothesis predicts the gap will largely
remain: kids will still be 10$\times$ steeper than LMs even when LM
training is CDS-matched. If the gap shrinks substantially, the
input-distribution explanation is doing more work than we currently
attribute to it.

\begin{thebibliography}{9}

\bibitem{chang2022}
Chang, T.\ A., \& Bergen, B.\ K.~(2022).
Word Acquisition in Neural Language Models.
\emph{Transactions of the Association for Computational Linguistics}, 10, 1--16.

\bibitem{daunizeau2017}
Daunizeau, J.~(2017).
Semi-analytical approximations to statistical moments of
sigmoid and softmax mappings of normal variables.
\emph{arXiv:1703.00091}.

\bibitem{feng2026}
Feng, S., et al.~(2026).
Child-directed-speech-trained language models. (Pending --- to be
filled in once we read the paper.)

\end{thebibliography}

\end{document}
